% JAC1964-20261121, /root/arithmetic_sign.  Fermat-specific tensor refinement.
\subsection*{AR-25. The Fermat gate as explicit tensor correspondences on a curve}

\begin{proposition}[AR-25]
Let $A$ be the Fermat Jacobian and $N=h^3(A)^G$ its exact
complementary motive.  For any $z\in\operatorname{CH}^2(N)_\mathbf Q$,
choose $W,C,f,M$ by the strengthened AR-21 construction with
$0\in W$.  After a finite extension choose a point $e\in C$ with
$f(e)=0$, and an isogeny
\[
 \phi:A\longrightarrow E^3,\qquad E:y^2=x^3-x,
 \qquad D=\deg\phi.
\]
Write $f_j=\operatorname{pr}_j\phi f:C\to E$ and let
$\psi=f_1\times f_2\times f_3:C^3\to E^3$, of degree
$\delta=\prod_j\deg f_j$.  Let $u=[i]$ on $E$, let
$\pi_1=\Delta_E-[0]\times E-E\times[0]$, and put
\[
 H=\{(a,b,c)\in(\mathbf Z/4)^3:a+b+c=0\},\qquad
 \Pi=\frac1{16}\sum_{(a,b,c)\in H}
    (u^a,u^b,u^c)_*\circ(\pi_1\otimes\pi_1\otimes\pi_1).
                                                               \tag{AR.25a}
\]
Then the following are exact Chow identities:
\[
 \Pi=\frac1D\phi_*\pi_N\phi^*,\qquad
 \Pi_*\psi_*\Delta_e=M\phi_*z.                         \tag{AR.25b}
\]
In particular
\[
 E_C=\frac1\delta\psi^*\Pi_*\psi_*
\]
is a self-transpose Chow idempotent on $C^3$, and
\[
 h_K(z,z)=\frac{\delta}{D M^2}
          h_{K'}((E_C)_*\Delta_e,(E_C)_*\Delta_e).       \tag{AR.25c}
\]
For $a\in\mathbf Z/4$ set
\[
 t_{j,a}=f_j^*\circ\pi_1\circ(u^a)_*\circ\pi_1\circ f_{j*}
       \in\operatorname{CH}^1(C\times C)_\mathbf Q.
\]
These are correspondences of degree zero in both factors,
normalized at the divisor $e$.  If $\widehat t_{j,a}$ denotes its
admissible metrized extension in Zhang's Theorem 2.3.5, then
\[
 h_K(z,z)=\frac1{16D M^2}
    \sum_{(a,b,c)\in H}
       \bigl\langle\widehat t_{1,a}\widehat t_{2,b}
                                  \widehat t_{3,c}\bigr\rangle_{K'}.
                                                               \tag{AR.25d}
\]
Here the arithmetic intersections are divided by $[K':\mathbf Q]$.
Thus the whole Fermat sign problem is a particular family of
explicit tensor-correspondence inequalities on curves.  Formula
(AR.25d) is an identity, not a proof that its right side is
nonnegative.
\end{proposition}

\begin{proof}
The generation assertion in AR-24 makes every $f_j$ nonconstant,
so each is a finite morphism of smooth curves.  Hence $\psi$ is
finite of degree $\delta>0$.  The chosen point satisfies $f(e)=0$.
Therefore
\[
 f_{j*}e=[0]\quad\hbox{in }\operatorname{CH}_0(E)_\mathbf Q
                                                               \tag{AR.25e}
\]
for every $j$, even as zero-cycles.  The extra point condition in
AR-21 leaves its Chow identity and denominator $M$ unchanged.
The finite extension only makes a point above the origin rational,
so it contributes no factor to normalized heights.  In particular
Zhang's tensor theorem is applied with an actual rational base
point, exactly as formulated in his \S2.2--2.3.

We check the first equality of (AR.25b) at the Chow level.  The
isogeny identities on rational Chow motives are
\[
 \phi_*\phi^*=D\,\mathrm{id},\qquad
 \phi^*\phi_*=D\,\mathrm{id}.                           \tag{AR.25f}
\]
The second uses the sum of translations by kernel points and the
fact that torsion translations act trivially on rational Chow
motives, as in AR-8; it is not an identity for an arbitrary
generically finite map.  Thus
$R=D^{-1}\phi_*\pi_N\phi^*$ is an exact self-transpose idempotent
supported on $h^3(E^3)$.

Let $a=[1+i]$ act diagonally on $E^3$ and write $U=a^*$ on
$h^3(E^3)$.  The exact endomorphism identity $a^4=[-4]$ and the
Chow--Kunneth multiplication eigenvalue give
$U^4=-64$.  The two coprime factors of $T^4+64$ are
\[
 F_{\rm ext}(T)=T^2+4T+8,\qquad
 F_{\rm lev}(T)=T^2-4T+8.
\]
Consequently
\[
 P_{\rm ext}=\frac{U^3-8U+32}{64}
\]
is an exact idempotent on $h^3(E^3)$: its polynomial is one modulo
$F_{\rm ext}$ and zero modulo $F_{\rm lev}$.  On realization it
selects precisely the types $(3,0),(0,3)$.

The projector $R$ commutes with $U$.  Indeed the corresponding
central rational CM endomorphism on $A$ commutes with every element
of $G$ and with $\pi_3$, so it commutes with $\pi_N=\pi_3p_G$;
the isogeny transports this central endomorphism to diagonal
$[1+i]$.  Also $\Pi$ commutes with $U$, directly from (AR.25a).
The group $H$ has order $16$.  On $H^1(E)^{\otimes3}$ its invariant
characters are exactly the all-holomorphic and all-antiholomorphic
ones: a character $(\epsilon_1,\epsilon_2,\epsilon_3)$ with
$\epsilon_j=\pm1$ is trivial on $H$ if and only if all three
$\epsilon_j$ agree.  Thus $R,\Pi$, and $P_{\rm ext}$ have the same
realization.  Since $R$ commutes with $P_{\rm ext}$,
$R(1-P_{\rm ext})$ and $P_{\rm ext}(1-R)$ are homologically
trivial idempotents.  Kimura nilpotence for the abelian motive makes
each zero.  Therefore $R=P_{\rm ext}$.  The same argument gives
$\Pi=P_{\rm ext}$, proving the asserted exact identification.
This argument uses the actual central polynomial and commuting
idempotents, rather than identifying Chow projectors merely from
their realizations.

The small diagonal pushes under $\psi$ to $\phi_*[W]$.  Every
other term of $\psi_*\Delta_e$ has at least one factor
$f_{j*}e=[0]$.  The factor $\pi_1$ kills this zero-cycle exactly,
so all six other terms vanish after applying $\Pi$.  Hence
\[
 \Pi_*\psi_*\Delta_e=\Pi_*\phi_*[W]
      =\phi_*(\pi_N)_*[W]=M\phi_*z.
\]
Here $\pi_N$ kills every Beauville grade except one and
$[W]_{(1)}=\operatorname{Ce}(W)=Mz$.

As in AR-24, graph composition and self-transposition prove that
$E_C$ is an exact self-transpose idempotent.  It sends $\Delta_e$
to $(M/\delta)\psi^*\phi_*z$.  Height adjointness and (AR.25f)
give
\[
 \begin{split}
 h_{K'}(E_C\Delta_e,E_C\Delta_e)
 &=\frac{M^2}{\delta}
       h_{K'}(\phi_*z,\phi_*z)\\
 &=\frac{D M^2}{\delta}h_K(z,z),
 \end{split}
\]
which proves (AR.25c).  These are raw middle height pairings of
codimension-two cycles on threefolds; no power of a polarization
occurs in their adjunction formula.  The factor $D$ explicitly
accounts for the isogeny, and no assertion that $\phi$ preserves
the principal polarization is required.  Self-transposition of
$R$ follows from graph transposition and the displayed $1/D$,
not from a polarization-preserving assumption on $\phi$.

Finally graph products give
\[
 E_C=\frac1{16\delta}\sum_{(a,b,c)\in H}
       t_{1,a}\otimes t_{2,b}\otimes t_{3,c}.            \tag{AR.25g}
\]
The factors $\pi_1$ make each $t_{j,a}$ have degree zero in
both variables.  Equation (AR.25e) gives
$(t_{j,a})_*e=(t_{j,a})^*e=0$, so these are precisely
normalized correspondences in Zhang's $\operatorname{Pic}^{-}$.
Apply his established Theorem 2.3.5 to each tensor term in
$h(\Delta_e,E_C\Delta_e)$, and use $E_C^2=E_C=E_C^t$.
The source writes the tensor action as pullback; transpose
adjointness and symmetry of the middle height give the same
pairing with $\Delta_e$ in both entries for the pushforward
convention here.
Together with (AR.25c), this yields (AR.25d).
\end{proof}

\paragraph{Precise relation to the cited positivity conjecture.}
The operator $E_C$ is a finite rational sum of tensors of normalized
curve correspondences.  Thus Zhang 2010, Conjecture 2.4.1, p.~28,
would imply $h(E_C\Delta_e,E_C\Delta_e)\ge0$ and hence the desired
semipositivity on the entire Fermat group, by (AR.25c).  Its
established Theorem 2.3.5, p.~27, proves only the triple-intersection
identity used here.  The source does not prove that conjectural
quadratic form is semipositive, and the divisor Hodge index theorem
of AR-22 does not determine the sign of its cubic intersections.

Moreover $E_C\Delta_e$ is nonzero exactly when $z$ is nonzero:
the isogeny induces an isomorphism on rational Chow groups and
$\psi^*$ is injective because $\psi_*\psi^*=\delta$.  Hence strict
positivity on this particular family of projected modified
diagonals is equivalent to the full strict Fermat gate.  No
equality characterization is added to Zhang's Conjecture 2.4.1,
whose displayed assertion is semipositivity.  The reduction
constructs the relevant tensors and normalizations; their sign is
the remaining problem.
